\section{Concavity of the Ising field functional}\label{sec:field}

The lower bound will identify a field whose conditional two-spin
products equal the overlap quantile. To turn that identity into the
supremum defining $S$, we need concavity of $\ell$ in its covariance
levels. We prove this property with the cascade weights held fixed.
We also establish continuity under refinement and describe exactly
which overlap quantiles have finite $S$.

\subsection{Conditional magnetizations and covariance derivatives}

Fix weights $w_i>0$, endpoints $\zeta_i$, and covariance levels
$0\leq h_0\leq\cdots\leq h_{k-1}$. Put
$v_i=h_i-h_{i-1}$, with $h_{-1}=0$, and let $\gamma_v$ denote the
centered Gaussian law of variance $v$, including the point mass at zero
when $v=0$. Recall the Gaussian operators $T_{\zeta,v}$ in the
definition of $\ell$. Let $F_{k-1}(z)=\log\cosh z$ and
$F_{i-1}=T_{\zeta_i,v_i}F_i$. Each $F_i$ is even, convex, smooth,
and 1-Lipschitz. Evenness follows by reflection; convexity follows from
H\"older's inequality for the logarithmic expectation, and ordinary
expectation when $\zeta=0$; the Lipschitz bound is preserved by either
operator.

The tilted position chain has $Z_{-1}=0$ and transitions
\begin{equation}\label{fld:kernel}
 \Prob(Z_i\in dy\mid Z_{i-1}=z)
 =e^{\zeta_i(F_i(y)-F_{i-1}(z))}\,\gamma_{v_i}(d(y-z)).
\end{equation}
At $i=0$ the exponential factor is one. Set $m_i=F_i'$ and
\begin{equation}\label{fld:B}
 B_i(h)=\E[m_i(Z_i)^2].
\end{equation}
The chain in \eqref{fld:kernel} is precisely the sampled-mark chain
of Proposition~\ref{arr:rpc} for the single-site Ising weight.

\begin{proposition}\label{fld:derivative}
The variables $m_i(Z_i)$ are the conditional means of
$\tanh Z_{k-1}$ along the tilted chain. In particular,
$0\leq B_0(h)\leq\cdots\leq B_{k-1}(h)\leq1$, and $B_i$ is the
expected product of two Ising spins whose fields have common level
$i$. For fixed weights, covariance differentiation gives
\begin{equation}\label{fld:gradient}
 d\ell(h)=-\frac12\sum_{i=0}^{k-1}w_iB_i(h)\,dh_i.
\end{equation}
The formula holds on the strict cone
$0<h_0<\cdots<h_{k-1}$ and, by continuity, for one-sided derivatives
along line segments in the closed cone. Subdividing an interval on
which $h$ is constant leaves $\ell$ unchanged and repeats its value of
$B$ on the resulting intervals.
\end{proposition}

\begin{proof}
Differentiating $T_{\zeta_i,v_i}F_i$ in its starting position gives
\[
 m_{i-1}(z)=\E[m_i(Z_i)\mid Z_{i-1}=z].
\]
At the terminal level $m_{k-1}=\tanh$. Thus these conditional means
form a martingale, and conditional Jensen proves the monotonicity of
their second moments. If two sampled paths separate immediately after
level $i$, their future marks and final spins are conditionally
independent given the common history. Their conditional product is
$m_i(Z_i)^2$. Proposition~\ref{arr:rpc} then proves the asserted
two-spin interpretation, including equal leaves at the last level.

Apply Lemma~\ref{arr:gaussian-calculus} to the single-site cascade
partition function with energy $z(\alpha)\sigma$, uniform prior on
$\sigma\in\{-1,1\}$, and covariance
$\E[z(\alpha)z(\alpha')]=h_{L(\alpha,\alpha')}$. The diagonal is
$h_{k-1}$. Its derivative cancels the derivative of the subtraction
$h_{k-1}/2$. The pair contribution at level $i$ is $B_i$, and that
level has probability $w_i$ after the tilt by
Proposition~\ref{arr:rpc}. This proves \eqref{fld:gradient}.

All tilted Gaussian integrals and their conditional means depend
continuously on the variances, also when some vanish. This follows by
Gaussian coupling and dominated convergence: the payoffs have at most
linear growth, their derivatives are bounded, and their exponential
weights have locally uniform Gaussian moments. This proves the
boundary statement. If an increment is zero, the adjacent positions
and backward functions coincide. The corresponding conditional means
and their squares therefore coincide as well, proving refinement
consistency.
\end{proof}

\subsection{Two monotonicity properties}

The variables $h_i$ are cumulative covariance levels. Increasing $h_j$
increases $v_j$ and, unless $j=k-1$, decreases $v_{j+1}$; thus the
mixed derivatives compare coupled changes of increments. The next two
comparisons control these derivatives and the response of $B_i$ to
scaling all covariance levels. Together they determine the sign of the
whole Hessian.

We first record a property of even Gaussian tilts. For $v>0$, an even
weight $e^{\zeta F(y)}$, and an absolute parent position $z\geq0$,
the density of the absolute child position $u\geq0$ is proportional to
\begin{equation}\label{fld:folded}
 e^{-u^2/(2v)}\cosh(zu/v)e^{\zeta F(u)}.
\end{equation}
For $z_2\geq z_1\geq0$, the likelihood ratio of the two densities is
increasing in $u$, since
\[
 \frac{d}{du}\log\frac{\cosh(z_2u/v)}{\cosh(z_1u/v)}
 =\frac{z_2\tanh(z_2u/v)-z_1\tanh(z_1u/v)}{v}\geq0.
\]
Consequently the absolute-position transition is stochastically
increasing in the absolute parent. Moreover, changing the even payoff
by a function nondecreasing, respectively nonincreasing, on
$[0,\infty)$ increases, respectively decreases, that transition in
stochastic order. These assertions follow directly by integrating an
increasing test against the ordered likelihood ratios.

\begin{lemma}\label{fld:mixed}
On the strict covariance cone, for every $i<j$,
\[
 \partial_{h_j}B_i(h)\leq0.
\]
\end{lemma}

\begin{proof}
All derivatives below exist by differentiating smooth Gaussian
convolutions; on compact subsets of the strict cone, linear growth of
the payoffs gives locally uniform bounds for this differentiation.
For $l<j$, covariance differentiation of the backward recursion gives
\begin{equation}\label{fld:backward-derivative}
 \partial_{h_j}F_l(z)
 =c_j-\frac{w_j}{2}
    \E[m_j(Z_j)^2\mid Z_l=z],
\end{equation}
where $c_j$ is independent of $z$ and vanishes unless $j=k-1$.
To verify the identity, increase the variance entering level $j$ and
decrease the variance entering level $j+1$. The two Gaussian generators
insert respectively
$(F_j''+\zeta_j m_j^2)/2$ and
$(F_j''+\zeta_{j+1}m_j^2)/2$, propagated by the same conditional
expectations. Their difference is $-w_jm_j^2/2$. At the last level,
$F_{k-1}''+m_{k-1}^2=1$ gives the additional constant $1/2$.
Propagation through the earlier operators gives
\eqref{fld:backward-derivative}.

The function $m_j$ is odd and nondecreasing, so $m_j^2$ is even and
nondecreasing in absolute position. By \eqref{fld:folded}, its
conditional expectation on the right-hand side of
\eqref{fld:backward-derivative} has the same property. Thus
$\partial_{h_j}F_l$ is nonincreasing on $[0,\infty)$, up to its
irrelevant constant. In particular, increasing $h_j$ decreases
$m_i(z)$ for each $z\geq0$ when $i<j$.

This change also decreases the sampled absolute positions up to
level $i$. Indeed, their Gaussian variances are unchanged, and the
change in every backward tilt $F_l$, $l\leq i$, is nonincreasing in
absolute position by the identity just proved. The likelihood-ratio
comparison following \eqref{fld:folded}, applied successively to the
transitions, gives this stochastic decrease. The root transition is
unchanged because $\zeta_0=0$. Combining the pointwise decrease of
$m_i(z)^2$ with the stochastic decrease of $|Z_i|$ proves the claim.
\end{proof}

The scaling comparison below is closely related to the monotonicity of
squared conditional magnetizations in Auffinger and Chen's proof of
raywise concavity \cite[Theorem~1 and its proof]{AuffingerChen2016}.

\begin{lemma}\label{fld:radial}
For each fixed monotone covariance vector $h$, $B_i(ch)$ is
nondecreasing in $c>0$. On the strict cone, therefore,
\begin{equation}\label{fld:radial-derivative}
 \sum_jh_j\partial_{h_j}B_i(h)\geq0.
\end{equation}
\end{lemma}

\begin{proof}
Use the original Gaussian increments of variances $v_i$, and replace
the terminal payoff by
$g_c(z)=\log\cosh(\sqrt c\,z)$. Write $F_i^{(c)}$ for these backward
values. This is the same recursion as at covariance $ch$, in the
original position units:
\[
 F_i^{(c)}(z)=F_i^{ch}(\sqrt c\,z).
\]
For $c_2\geq c_1>0$, the difference $g_{c_2}-g_{c_1}$ is even and
nondecreasing on $[0,\infty)$. This property propagates through every
backward operator. For a positive exponent, the difference of the
two backward values can be written
\[
 \frac1\zeta\log\E_{c_1,z}
          e^{\zeta(F^{(c_2)}-F^{(c_1)})(Y)},
\]
where the expectation is in the transition tilted by the lower-$c$
payoff. For exponent zero it is the ordinary expectation of the
difference. In either case \eqref{fld:folded} shows that the result
is nondecreasing in the absolute starting position. It is even by
reflection. Hence the tilted absolute-position chain is stochastically
increasing with $c$.

It remains to compare conditional spin means, rather than derivatives
in the rescaled position. Define
\[
 a_i^{(c)}(z)=
 \E_c[\tanh(\sqrt c\,Z_{k-1})\mid Z_i=z].
\]
We prove backwards that these functions are odd and, on $z\geq0$,
nonnegative and nondecreasing both in $z$ and in $c$. This is true at
the terminal level. Conditional on the absolute value $u$ of a child
with parent $z\geq0$, the mean sign of the child is
$\tanh(zu/v)$: the even tilt cancels in the ratio of its positive and
negative densities. The backward conditional mean is thus the integral
of
\[
 a_{i+1}^{(c)}(u)\tanh(zu/v)
\]
against the folded tilted transition. Its factors are nonnegative and
nondecreasing in $u$, while the transition is increasing in $z$ and
$c$. The induction follows, with zero increments obtained by
continuity.

Finally, $a_i^{(c)}(z)=m_i^{ch}(\sqrt c\,z)$. These are precisely
the physical conditional spin means, without an additional scaling
factor. Combining their pointwise increase with that of the sampled
absolute positions shows that
$\E_c[a_i^{(c)}(Z_i)^2]=B_i(ch)$ increases with $c$. Differentiating
at $c=1$ gives \eqref{fld:radial-derivative}.
\end{proof}

\subsection{Concavity, continuity, and the domain of \texorpdfstring{$S$}{S}}

\begin{theorem}\label{fld:concavity}
For fixed interval weights, $\ell$ is concave on the closed cone of
nonnegative nondecreasing covariance levels. For any two such step
paths $h,v$, represented on a common subdivision,
\begin{equation}\label{fld:tangent}
 \ell(v)\leq\ell(h)
   -\frac12\int_0^1B_h(s)(v(s)-h(s))\,ds,
\end{equation}
where $B_h$ is the step path with values $B_i(h)$. Furthermore,
\begin{equation}\label{fld:lipschitz}
 |\ell(h)-\ell(v)|\leq\frac12\|h-v\|_{L^1(0,1)}.
\end{equation}
Thus $\ell$ has a unique Lipschitz extension to bounded nonnegative
nondecreasing paths, independent of the step approximations used.
\end{theorem}

\begin{proof}
On the strict cone put
$M_{ij}=w_i\partial_{h_j}B_i$. Equation~\eqref{fld:gradient} shows
that $M$ is symmetric and equals $-2$ times the Hessian of $\ell$.
Lemma~\ref{fld:mixed} and symmetry give $M_{ij}\leq0$ for $i\ne j$.
Let $D=\diag(h_0,\ldots,h_{k-1})$. By
Lemma~\ref{fld:radial}, every row sum of $DMD$ is nonnegative:
\[
 \sum_j(DMD)_{ij}
 =h_iw_i\sum_jh_j\partial_{h_j}B_i\geq0.
\]
If a symmetric matrix $A$ has nonpositive off-diagonal entries and
row sums $r_i\geq0$, then
\[
 x^{\mathsf T}Ax
 =\sum_i r_ix_i^2+
    \sum_{i<j}(-A_{ij})(x_i-x_j)^2\geq0.
\]
Applying this identity to $DMD$, and using invertibility of $D$ on the
strict cone, proves that $M$ is positive semidefinite. Hence the Hessian
of $\ell$ is negative semidefinite.

Integrating along line segments gives concavity and the supporting
inequality in the strict cone. Approximate boundary covariance vectors
by strict ones and use the continuity in
Proposition~\ref{fld:derivative} to obtain
\eqref{fld:tangent} on the closed cone. If $h$ and $v$ initially use
different intervals, put both on their common subdivision. Refinement
consistency makes this operation harmless. Finally, integrate
\eqref{fld:gradient} along the segment from $h$ to $v$ and use
$0\leq B_i\leq1$ to prove \eqref{fld:lipschitz}. Monotone step paths
are dense in the bounded monotone paths in $L^1$, proving the extension.
\end{proof}

The sign in \eqref{fld:tangent} is the one needed in the cavity
argument: if a limiting field satisfies $B_h=p$, that field maximizes
$\ell(v)+\frac12\int pv$ over admissible $v$.

In Ho's normalization \cite[Theorem~1.8]{Ho2026},
$\ell(h)=-\psi(h/2)$. Thus his covariance-path theorem gives the
zero-field concavity above. Both conventions use the terminal function
$\log\cosh$, so no additive $\log2$ term occurs. The preceding proof
also gives the supporting inequality and endpoint statements in the form
used below.

The functional $S$ is finite except at the single extremal quantile
$p=1$. A quantitative bound will also clarify the trial domain in the
upper comparison.

\begin{lemma}\label{fld:domain}
For $p\in\cP$, put $d=\int_0^1(1-p(s))\,ds$. If $d>0$, then
\begin{equation}\label{fld:domain-bound}
 0\leq S(p)\leq(d^{-1}-1)\log2.
\end{equation}
If $d=0$, so that $p=1$ almost everywhere, then $S(p)=+\infty$.
Moreover $S$ is increasing in the pointwise order on quantiles.
\end{lemma}

\begin{proof}
The zero field gives $S(p)\geq0$, and monotonicity follows because
every admissible field is nonnegative. Suppose $0<d\leq1$, and define
\[
 g_d(z)=\frac1d\log(2\cosh(dz))-\log2.
\]
Then $g_d\geq\log\cosh$: with $r=1/d\geq1$, the inequality
$(a+b)^r\geq a^r+b^r$ applied to $a=e^{dz}$ and $b=e^{-dz}$ gives
this assertion. For $0\leq\zeta\leq1$ and $v\geq0$,
\begin{equation}\label{fld:majorant}
 T_{\zeta,v}g_d\leq g_d+\frac12\max(\zeta,d)v.
\end{equation}
Indeed, if $\zeta\geq d$, apply Minkowski's inequality with exponent
$r=\zeta/d\geq1$ to
$e^{dz+dG}+e^{-dz-dG}$, where $G\sim\gamma_v$.
The $L^r$ norm of either exponential gains the factor
$e^{rd^2v/2}$, proving the bound $\zeta v/2$. If
$0<\zeta\leq d$, Jensen's inequality for the concave power
$\zeta/d$ gives the bound $dv/2$. At $\zeta=0$, apply Jensen directly
to the logarithm to obtain the same bound.

Iterate \eqref{fld:majorant} through the field recursion. Since
$g_d(0)=(d^{-1}-1)\log2$, this yields
\[
 \ell(h)\leq(d^{-1}-1)\log2
 -\frac12\sum_i(1-\max(\zeta_i,d))(h_i-h_{i-1})
 =(d^{-1}-1)\log2-\frac12\int_d^1h(s)\,ds.
\]
For any nonnegative nondecreasing $h$,
\[
 \int_0^1(1-p(s))h(s)\,ds\geq\int_0^d h(s)\,ds.
\]
To see this, the density $1-p$ lies in $[0,1]$ and has mass $d$;
among all such densities the integral against an increasing function
is minimized by $\ind_{(0,d)}$. For $d=1$ the assertion is immediate.
For $0<d<1$, subtract this indicator and then subtract the constant
$h(d)$ from the integrand: the resulting product is nonnegative on both
sides of $d$. Combining the last two
displays gives \eqref{fld:domain-bound} after taking the supremum over
$h$.

Finally, for $p=1$ take the constant field $h(s)=H$. Its recursion is
ordinary expectation at the root, and
\[
 \ell(h)+\frac12\int_0^1h(s)\,ds
 =\E\log\cosh(\sqrt H\,G)
 \geq\sqrt H\,\E|G|-\log2\longrightarrow\infty,
\]
where $G$ is standard Gaussian. This proves the last case.
\end{proof}
